%!TEX root = ../template.tex
%%%%%%%%%%%%%%%%%%%%%%%%%%%%%%%%%%%%%%%%%%%%%%%%%%%%%%%%%%%%%%%%%%%%
%% options-ext.tex
%% NOVA thesis documentation file
%%
%% Detailed explanation of options-ext.sty functionality.
%%%%%%%%%%%%%%%%%%%%%%%%%%%%%%%%%%%%%%%%%%%%%%%%%%%%%%%%%%%%%%%%%%%%

\typeout{NT FILE options-ext.tex}

\chapter{Explanation of options-ext.sty}
\label{annex:options-ext}

The file \texttt{NOVAthesisFiles/StyFiles/options-ext.sty} provides a set of powerful extensions to the \texttt{options} package (originally by Daan Leijen). These extensions are crucial for the complex key-value management system that \emph{NOVAthesis} uses to handle multi-school, multi-language configurations.

\section{Choice Iteration}
The command \texttt{\textbackslash optionchoicedo\{option\}\{code\}} is introduced to iterate over all possible values of a ``choice'' type option. This allows the template to perform actions for every school or language defined in the system.

\section{Custom Key Handlers}
This package adds several new ``handlers'' (special actions that can be triggered when a key is set) to the \texttt{options} framework:

\begin{itemize}
    \item \textbf{Copying Handlers}: 
        \begin{itemize}
            \item \texttt{.copy list}: Allows copying an entire list from one key to another.
            \item \texttt{.copy choice}: Copies choice definitions from a source key to a target.
        \end{itemize}
    \item \textbf{Pushing and Merging}:
        \begin{itemize}
            \item \texttt{.newpush}: Automatically creates a list option if it doesn't exist and then pushes a value to it.
            \item \texttt{.push ifnew}: Pushes a value to a list only if it's not already there (ensuring uniqueness).
            \item \texttt{.newpush ifnew}: A combination of the above, creating the list if needed and then pushing only unique values.
        \end{itemize}
    \item \textbf{Utility Handlers}:
        \begin{itemize}
            \item \texttt{.size}: Counts the number of items in a list or the number of choices in a choice option and stores the result in a macro.
            \item \texttt{.show items}: A debugging tool that prints the current contents of a list option.
        \end{itemize}
\end{itemize}

\section{Significance}
These extensions allow the template's configuration files to be much more expressive and declarative. Instead of complex \texttt{\textbackslash if} logic, the template can manage schools, fonts, and styles by dynamically building and merging lists of options.

\section{Summary}
In summary, \texttt{options-ext.sty} is the engine that enables \emph{NOVAthesis} to handle its highly modular and customizable architecture through advanced key-value manipulation.
